% GENERATED FILE. Edit canonical lessons, then run npm run generate:books.
% canonical-source-hash: fnv1a64:26f5756e5d5228c8
% canonical-lessons: ZH-C482-jizhen, ZH-C482-zhiliao, ZH-C482-yaopian, ZH-C482-xiaodu, ZH-C482-huli

\chapter{Treatment}
\label{ch:zh-treatment}
\hlchaptermodality{482}

\begin{chapteropening}
\textbf{By the end of this chapter:} \emph{I can talk about the emergency department, treatment, tablets, cleaning a wound and nursing the sick.}

Complete the last lesson of chapter 482: I can talk about the emergency department, treatment, tablets, cleaning a wound and nursing the sick.
\end{chapteropening}

\section[jízhěn]{jízhěn — emergency department; urgent treatment}

\label{lesson:ZH-C482-jizhen}



\begin{quote}
Before the new one: say the Mandarin for to become infected, then the Mandarin for a virus.
\end{quote}

\subsection*{You'll want to know: jízhěn}
\emph{jízhěn} — \textquotedblleft{}emergency department; urgent treatment\textquotedblright{}.

\begin{quote}
\emph{wǒ xiǎng dāng jízhěn yīshēng} — I want to be an emergency doctor.
\end{quote}

Its characters come later; for now it is a word for the ear.

\begin{cousinweb}
\emph{jízhěn} is built from \emph{jí} 'urgent' and \emph{zhěn} 'to examine', both inherited from Middle Chinese, \emph{kip} and \emph{tsyinX} (Guangyun, in Baxter's transcription), which glosses the first 'urgent' and the second as feeling the pulse.
\end{cousinweb}

\begin{grammarlens}[title={What you've built}]
One more word for saying how it was.
\end{grammarlens}

\subsection*{Your turn}
\begin{itemize}
\raggedright
  \item \emph{Say it:} \emph{jízhěn}
  \item \emph{Say it:} \emph{jízhěn}, once more
  \item \emph{Say it:} \emph{wǒ xiǎng dāng jízhěn yīshēng} — I want to be an emergency doctor.
  \item \emph{Recall:} say the Mandarin for to become infected, then the Mandarin for a virus, then say \emph{jízhěn} again
\end{itemize}

\begin{tcolorbox}[breakable,colback=teal!4,colframe=teal!35!black,title={Before you move on}]
What does jízhěn mean? (\textquotedblleft{}Emergency department; urgent treatment\textquotedblright{}.) Say it once more.
\end{tcolorbox}
\section[zhìliáo]{zhìliáo — to give medical treatment; to cure}

\label{lesson:ZH-C482-zhiliao}



\begin{quote}
Before the new one: say the Mandarin for a virus, then the Mandarin for the emergency department.
\end{quote}

\subsection*{You'll want to know: zhìliáo}
\emph{zhìliáo} — \textquotedblleft{}to give medical treatment; to cure\textquotedblright{}.

\begin{quote}
\emph{tā zhìliáo le bàn nián xiànzài yǐjīng hǎo le} — He was treated for six months, and now he's well again.
\end{quote}

Its characters come later; for now it is a word for the ear.

\begin{cousinweb}
\emph{zhìliáo} is built from \emph{zhì} 'to put in order' and \emph{liáo} 'to heal', both inherited from Middle Chinese, \emph{dri} and \emph{ljewH} (Guangyun, in Baxter's transcription). The dictionary lists the first under three readings and glosses them 'to put in order'; it glosses the second 'to treat an illness'. The falling tone of \emph{zhì} follows the reading \emph{driH}.
\end{cousinweb}

\begin{grammarlens}[title={What you've built}]
One more word for saying how it was.
\end{grammarlens}

\subsection*{Your turn}
\begin{itemize}
\raggedright
  \item \emph{Say it:} \emph{zhìliáo}
  \item \emph{Say it:} \emph{zhìliáo}, once more
  \item \emph{Say it:} \emph{tā zhìliáo le bàn nián xiànzài yǐjīng hǎo le} — He was treated for six months, and now he's well again.
  \item \emph{Recall:} say the Mandarin for a virus, then the Mandarin for the emergency department, then say \emph{zhìliáo} again
\end{itemize}

\begin{tcolorbox}[breakable,colback=teal!4,colframe=teal!35!black,title={Before you move on}]
What does zhìliáo mean? (\textquotedblleft{}To give medical treatment; to cure\textquotedblright{}.) Say it once more.
\end{tcolorbox}
\section[yàopiàn]{yàopiàn — tablet; pill}

\label{lesson:ZH-C482-yaopian}



\begin{quote}
Before the new one: say the Mandarin for the emergency department, then the Mandarin for to give medical treatment.
\end{quote}

\subsection*{You'll want to know: yàopiàn}
\emph{yàopiàn} — \textquotedblleft{}tablet; pill\textquotedblright{}.

\begin{quote}
\emph{yīshēng shuō yàopiàn yì tiān chī sān cì} — The doctor says to take the pills three times a day.
\end{quote}

Its characters come later; for now it is a word for the ear.

\begin{cousinweb}
\emph{yàopiàn} is built from \emph{yào} 'medicine' and \emph{piàn} 'a slice', both inherited from Middle Chinese, \emph{yak} and \emph{phenH} (Guangyun, in Baxter's transcription), which explains the first, after an older book, as a plant that cures illness, and glosses the second 'half' and 'a split piece of wood'.
\end{cousinweb}

\begin{grammarlens}[title={What you've built}]
One more word for saying how it was.
\end{grammarlens}

\subsection*{Your turn}
\begin{itemize}
\raggedright
  \item \emph{Say it:} \emph{yàopiàn}
  \item \emph{Say it:} \emph{yàopiàn}, once more
  \item \emph{Say it:} \emph{yīshēng shuō yàopiàn yì tiān chī sān cì} — The doctor says to take the pills three times a day.
  \item \emph{Recall:} say the Mandarin for the emergency department, then the Mandarin for to give medical treatment, then say \emph{yàopiàn} again
\end{itemize}

\begin{tcolorbox}[breakable,colback=teal!4,colframe=teal!35!black,title={Before you move on}]
What does yàopiàn mean? (\textquotedblleft{}Tablet; pill\textquotedblright{}.) Say it once more.
\end{tcolorbox}
\section[xiāodú]{xiāodú — to disinfect}

\label{lesson:ZH-C482-xiaodu}



\begin{quote}
Before the new one: say the Mandarin for to give medical treatment, then the Mandarin for a pill.
\end{quote}

\subsection*{You'll want to know: xiāodú}
\emph{xiāodú} — \textquotedblleft{}to disinfect\textquotedblright{}.

\begin{quote}
\emph{hùshi xiān gěi shāngkǒu xiāodú} — The nurse disinfects the wound first.
\end{quote}

Its characters come later; for now it is a word for the ear.

\begin{cousinweb}
\emph{xiāodú} is built from \emph{xiāo} 'to wipe out' and \emph{dú} 'poison', both inherited from Middle Chinese, \emph{sjew} and \emph{dowk} (Guangyun, in Baxter's transcription), which glosses the first 'to destroy' and 'to use up' and the second 'harm'.
\end{cousinweb}

\begin{grammarlens}[title={What you've built}]
One more word for saying how it was.
\end{grammarlens}

\subsection*{Your turn}
\begin{itemize}
\raggedright
  \item \emph{Say it:} \emph{xiāodú}
  \item \emph{Say it:} \emph{xiāodú}, once more
  \item \emph{Say it:} \emph{hùshi xiān gěi shāngkǒu xiāodú} — The nurse disinfects the wound first.
  \item \emph{Recall:} say the Mandarin for to give medical treatment, then the Mandarin for a pill, then say \emph{xiāodú} again
\end{itemize}

\begin{tcolorbox}[breakable,colback=teal!4,colframe=teal!35!black,title={Before you move on}]
What does xiāodú mean? (\textquotedblleft{}To disinfect\textquotedblright{}.) Say it once more.
\end{tcolorbox}
\section[\texorpdfstring{\zh{护理} (hùlǐ)}{护理 (hùlǐ)}]{\zh{护理} (hùlǐ) — to care for the sick; to tend}

\label{lesson:ZH-C482-huli}



\begin{quote}
Before the new one: say the Mandarin for a pill, then the Mandarin for to disinfect.
\end{quote}

\subsection*{You'll want to know: \zh{护理}}
\textbf{\zh{护理}} — \emph{hùlǐ} — \textquotedblleft{}to care for the sick; to tend\textquotedblright{}.

\begin{quote}
\textbf{\zh{护理病人的时候你不能怕血}} \emph{hùlǐ bìngrén de shíhou nǐ bù néng pà xiě} — When you look after patients, you can't be afraid of blood.
\end{quote}

\begin{cousinweb}
\textbf{\zh{护理}} (\emph{hùlǐ}) is built from \zh{护} 'to protect' and \zh{理} 'to put in order', both inherited from Middle Chinese, \emph{huH} and \emph{liX} (Guangyun, in Baxter's transcription), which glosses the first 'to rescue' and 'to help' and the second 'to manage'; the sense of caring for the sick is not given there.
\end{cousinweb}

\begin{grammarlens}[title={What you've built}]
One more word for saying how it was.
\end{grammarlens}

\subsection*{Your turn}
\begin{itemize}
\raggedright
  \item \emph{Say it:} \emph{hùlǐ}
  \item \emph{Say it:} \emph{hùlǐ}, once more
  \item \emph{Say it:} \emph{hùlǐ bìngrén de shíhou nǐ bù néng pà xiě} — When you look after patients, you can't be afraid of blood.
  \item \emph{Recall:} say the Mandarin for a pill, then the Mandarin for to disinfect, then say \emph{hùlǐ} again
\end{itemize}

\begin{tcolorbox}[breakable,colback=teal!4,colframe=teal!35!black,title={Before you move on}]
What does \zh{护理} mean? (\textquotedblleft{}To care for the sick; to tend\textquotedblright{}.) Say it once more.
\end{tcolorbox}
